\section{The actual local deformation presentation and its relation count}
\label{sec:arithmetic-presentations}

\paragraph{The local ring that will be presented.}
Fix the original complete coefficient DVR \(O\), its finite residue
field \(k\) of characteristic \(\ell\), a frame
\(F.\mathrm{Ok}(C_0)\), and one supplied
\(S:\mathrm{Space}(O,F,n)\), where \(n=2\) or \(3\).
Write \(\mathcal A=S.A\), \(\mathcal B=S.B\), and
\(u=S.u\).  For a supplied original local object
\(L:S.\mathrm{Loc}(x)\), put
\[
 O'=L.O',\qquad R=L.R,\qquad
 \iota=L.\iota:\mathcal A\longrightarrow R,
 \qquad \rho=\iota_*u.                                      \tag{APD.1}
\]
All constructions below use these rings and this map.  In particular,
the final assertion is the literal presentation
\[
 R\simeq O'[[X_1,\ldots,X_N]]/(f_1,\ldots,f_{N'}),
 \qquad N=N'+nF.d+1.                                      \tag{APD.2}
\]
It holds for every supplied \(L\); it is not obtained by choosing a
more favorable point or a further coefficient extension.

The argument has three parts.  First we prove the commutative algebra
that converts an obstruction space into a bound on relations.  Next
we show that the actual local normalized representation represents its
deformation problem, including points with triangular, rather than
diagonal, reduction.  Finally we construct its continuous obstruction
cocycles, normalize their corrections, and compute the difference
between tangent and obstruction dimensions.  The frame's original
Euler field is used exactly as a field of \(F.\mathrm{Ok}(C_0)\).
Its arithmetic construction was supplied by LFC20; the original
\(P\) hypothesis at the arithmetic frame entry is not changed here.

\paragraph{Conventions and available inputs.}
For any local ring \(\mathcal O\), an
\(\mathrm{IsCO}(\mathcal O,C)\)-algebra is a complete Noetherian
local \(\mathcal O\)-algebra whose specified coefficient map induces
an isomorphism of residue fields.  Its residue map to
\(\kappa=\mathcal O/\mathfrak m_{\mathcal O}\) is denoted
\(\operatorname{res}_C\).  Every homomorphism between these coefficient
algebras is local, preserves this residue map, and is adically
continuous: the unit argument in ALC1 proves the first two facts,
and preservation of powers proves the third.  The general statements
below explicitly say when \(\mathcal O\) must itself be Noetherian
or complete; the arithmetic DVR assumptions are not imposed on the
broader statements.

We use the complete Noetherian algebra and Artin--Rees proofs in
ALC5--ALC7, the actual local chart and representation in ALC11 and
ALC17, the original weighted entry algebra in RTC7, its fixed-weight
continuity in RTC9, and the ordered flag and characteristic-polynomial
arguments in ACP3--ACP4.  We also use the elementary Nakayama lemma,
the Noetherian formal-power-series theorem, and compactness of an
inverse limit of finite discrete sets.  The constructions and all
deformation-specific applications of these facts are proved here.

\subsection{Minimal generators and obstruction spaces}

% BEGIN APD1
\paragraph{APD1. Tangent vectors, relative cotangents, and dual numbers.}
\label{claim:APD1}
Let \(\mathcal O\) be any commutative local ring and let
\(R\) be an \(\mathrm{IsCO}(\mathcal O,R)\)-algebra.  Put
\[
 T=\operatorname{Der}_{\mathcal O}(R,\kappa),\qquad
 V=\mathfrak m_R/(\mathfrak m_R^2+
                         \mathfrak m_{\mathcal O}R).
                                                               \tag{APD.3}
\]
Here the target of a derivation is an \(R\)-module by
\(\operatorname{res}_R\).  Restriction gives
\(T\simeq\operatorname{Hom}_\kappa(V,\kappa)\), so \(T\)
is finite dimensional.  There is a bijection
\[
 \operatorname{Hom}_{\mathcal O\text{-alg}}
       (R,\kappa[\epsilon]/\epsilon^2)\simeq T,
 \qquad f_D(r)=\operatorname{res}_R(r)+\epsilon D(r).
                                                               \tag{APD.4}
\]
This proves the full original \texttt{pr\_tan}, including its
cardinality assertion when the residue field is not finite.

\begin{proof}
A derivation kills \(\mathcal O\), and the Leibniz rule makes it
zero on \(\mathfrak m_R^2\) and on
\(\mathfrak m_{\mathcal O}R\).  Its restriction therefore gives
the indicated linear functional.  Conversely let \(a\) be such
a functional.  For \(r\in R\), choose \(c\in\mathcal O\)
with \(r-c\in\mathfrak m_R\), and set
\(D_a(r)=a([r-c])\).  A different choice changes \(r-c\) by an
element of \(\mathfrak m_{\mathcal O}R\), so the definition is
independent of the choice.  Choosing lifts for two residues shows
additivity and \(\mathcal O\)-linearity.  Expanding
\((c+y)(d+z)\), with \(y,z\in\mathfrak m_R\), shows
\(D_a(rs)=\bar rD_a(s)+\bar sD_a(r)\), since \(yz\) dies in
\(V\).  These two constructions are inverse.

The ideal \(\mathfrak m_R\) is finitely generated because \(R\)
is Noetherian.  Its images generate \(V\) over the residue field;
thus \(V\), and hence its dual, are finite dimensional.  No
Noetherian hypothesis on \(\mathcal O\) was used.

The residue field and the dual-number ring are both
\(\mathrm{IsCO}(\mathcal O,-)\)-algebras.  For the latter the
maximal ideal is \((\epsilon)\), its square is zero, and every
ideal is a finite-dimensional \(\kappa\)-space.  Thus any algebra
map to it has first coordinate \(\operatorname{res}_R\).
Its second coordinate is a derivation, by multiplying dual numbers.
Conversely the Leibniz rule proves that \(f_D\) in (APD.4) is
an algebra homomorphism.  This gives the claimed bijection, not
merely equality of two finite counts.
\end{proof}
% END APD1

% BEGIN APD2
\paragraph{APD2. The minimal power-series presentation is constructed.}
\label{claim:APD2}
Suppose in addition that \(\mathcal O\) is complete and
Noetherian.  If \(t=\dim_\kappa T\), there is a surjective
continuous \(\mathcal O\)-algebra map
\[
 p:P=\mathcal O[[X_1,\ldots,X_t]]\longrightarrow R,
 \qquad I=\ker p\subset
        \mathfrak m_{\mathcal O}P+\mathfrak m_P^2.             \tag{APD.5}
\]
Equivalently, every linear coefficient of every element of \(I\)
belongs to \(\mathfrak m_{\mathcal O}\).  The number of variables
is exactly the tangent dimension.

\begin{proof}
Choose \(z_1,\ldots,z_t\in\mathfrak m_R\) lifting a basis
of \(V\) in (APD.3).  Nakayama applied to the finite module
\(\mathfrak m_R/(\mathfrak m_{\mathcal O}R+(z_1,\ldots,z_t))\)
gives
\(\mathfrak m_R=\mathfrak m_{\mathcal O}R+(z_1,\ldots,z_t)\).
Evaluation of a power series at the \(z_i\) converges in \(R\)
and defines \(p\).  To prove surjectivity, first lift the residue
of a prescribed \(r\in R\) by an element of \(\mathcal O\).
At stage \(a\), an error in \(\mathfrak m_R^a\) is a sum of
products of \(a\) of the indicated maximal-ideal generators.
Replace the remaining coefficients by lifts of their residues from
\(\mathcal O\).  This lifts the error modulo
\(\mathfrak m_R^{a+1}\) by an element of \(\mathfrak m_P^a\).
Successive corrections converge in \(P\) and their image is \(r\).

If \(f\in I\), write its constant and linear terms as
\(c+\sum_i a_iX_i\).  The constant coefficient \(c\) lies
in \(\mathfrak m_{\mathcal O}\), because (p) preserves
residues.  Reducing \(p(f)=0\) in \(V\) gives
\(\sum_i\overline{a_i}[z_i]=0\).  The chosen basis forces
\(a_i\in\mathfrak m_{\mathcal O}\) for every \(i\), proving
(APD.5).  Conversely this coefficient condition is exactly the
displayed ideal containment, by separating terms of degrees zero,
one, and at least two.

The same approximation argument gives a useful criterion used below:
if \(f:R_1\to R_2\) is a local map of complete Noetherian local
rings, induces a surjection on residue fields, and
\(\mathfrak m_{R_2}=f(\mathfrak m_{R_1})R_2\), then \(f\)
is surjective.  A correction of order \(a\) is a sum of degree
\(a\) products of images; its coefficients can be lifted modulo
the maximal ideal, and completeness sums the resulting corrections.
This criterion does not require an injection of coefficient fields.
\end{proof}
% END APD2

% BEGIN APD3
\paragraph{APD3. The space of necessary relations and its coordinates.}
\label{claim:APD3}
For (APD.5), set \(C=P/\mathfrak m_PI\) and
\(K=I/\mathfrak m_PI\).  Then \(C\to R\) is a surjective
map of \(\mathrm{IsCO}(\mathcal O,-)\)-algebras,
\(\mathfrak m_CK=0\), and \(K\) is a finite-dimensional
\(\kappa\)-space.  A basis of \(K\) lifts to a generating
family of the ideal \(I\).

More generally, let \(C\) be any such coefficient algebra and
let \(K\subset C\) be any ideal annihilated by
\(\mathfrak m_C\).  Suppose a finite family \((e_i)\)
generates \(K\), and every relation
\(\sum_i a_ie_i=0\), with \(a_i\in\mathcal O\), has all
\(a_i\in\mathfrak m_{\mathcal O}\).  Then there are
\(\mathcal O\)-linear coordinates \(c_i:K\to\kappa\) such
that every \(y\in K\) can be written
\[
 y=\sum_i a_ie_i,\qquad \overline{a_i}=c_i(y).
                                                               \tag{APD.6}
\]
This is the full original \texttt{pr\_dualb}.

\begin{proof}
The finite-variable power-series ring \(P\) is complete
Noetherian local.  Its ideals are closed, and its quotients are
complete, by exactness of finite-module completion in ALC7.
Both quotient rings therefore have the asserted coefficient
structures.  Multiplication by \(\mathfrak m_C\) on \(K\)
is zero by construction.  Since \(I\) is a finite \(P\)-module,
\(I/\mathfrak m_PI\) is finite dimensional.  Nakayama applied
to \(I\) shows that lifts of any basis generate \(I\).

In the more general assertion multiplication by an element of \(C\)
on \(K\) depends only on its residue, so \(K\) is a
\(\kappa\)-space.  Any ideal expression for \(y\) in terms
of the \(e_i\) can replace its coefficients in \(C\) by
residue lifts from \(\mathcal O\).  Thus the \(e_i\) span
over \(\kappa\).  A linear dependence over \(\kappa\)
can also be lifted to \(\mathcal O\); its value is zero because
\(\mathfrak m_CK=0\).  The stated relation hypothesis makes
every coefficient zero.  Hence \((e_i)\) is a basis.
Take its coordinate functionals and lift their values back to
\(\mathcal O\); this gives (APD.6) and proves linearity as well.
The empty-family case gives \(K=0\) and is included.
\end{proof}
% END APD3

% BEGIN APD4
\paragraph{APD4. A nonzero minimal relation cannot split off.}
\label{claim:APD4}
Let \(C\twoheadrightarrow R\) be the small extension of APD3,
and let \(u:K\to\kappa\) be a nonzero linear functional.
The quotient
\[
 C_u=C/\ker u\longrightarrow R                            \tag{APD.7}
\]
does not admit a \(\mathcal O\)-algebra section.

\begin{proof}
Every \(\kappa\)-subspace of \(K\) is an ideal of \(C\),
because its scalar action factors through \(C/\mathfrak m_C\).
Thus \(C_u\) is a complete Noetherian local coefficient algebra,
and its kernel over \(R\) is the nonzero one-dimensional space
\(K/\ker u\).  Write this kernel as \(J\).
Minimality (APD.5), after passing to \(C_u\), gives
\[
 J\subset\mathfrak m_{\mathcal O}C_u+\mathfrak m_{C_u}^2.
                                                               \tag{APD.8}
\]
Suppose \(s:R\to C_u\) were a section.  Every
\(a\in\mathfrak m_{C_u}\) differs from \(s(p(a))\) by an
element of \(J\).  The first summand on the right of (APD.8)
already belongs to \(s(\mathfrak m_R)C_u\).  Consequently
\[
 \mathfrak m_{C_u}=s(\mathfrak m_R)C_u+\mathfrak m_{C_u}^2.
\]
Nakayama makes the last square unnecessary.  The surjectivity
criterion proved in APD2 now makes \(s\) surjective.  Since
\(p\circ s=1\), a surjective section makes \(p\) an
isomorphism, contradicting \(J\ne0\).  This argument proves
the obstruction to splitting, rather than assuming a minimal
presentation has a universal obstruction theory.
\end{proof}
% END APD4

% BEGIN APD5
\paragraph{APD5. Full obstruction-to-presentation algebra.}
\label{claim:APD5}
Let \(\mathcal O\) be complete Noetherian local with finite
residue field \(\kappa\), let \(R\) be an
\(\mathrm{IsCO}(\mathcal O,R)\)-algebra, and let \(H\) be
a finite-dimensional \(\kappa\)-space.  Suppose the original
\(\mathrm{HasObs}(\mathcal O,R,H)\) holds, in the following
literal sense.  For every surjection \(p:C\to R\) of coefficient
algebras with \(\mathfrak m_C\ker p=0\), there is a linear map
\[
 \mathrm{ob}_p:\operatorname{Hom}_{\mathcal O}(\ker p,\kappa)
                      \longrightarrow H                 \tag{APD.9}
\]
such that, whenever \(\mathrm{ob}_p(u)=0\), every factorization
\(C\overset q\twoheadrightarrow C'\overset{p'}\longrightarrow R\)
with \(p'q=p\) and
\(u(y)=0\Rightarrow q(y)=0\) for \(y\in\ker p\)
has an \(\mathcal O\)-algebra section of \(p'\).
Then there is a surjection
\[
 \mathcal O[[X_1,\ldots,X_t]]\longrightarrow R,
 \qquad t=\dim_\kappa T,
                                                               \tag{APD.10}
\]
whose kernel is generated by a specified family of exactly
\(h=\dim_\kappa H\) elements.  This proves the full original
\texttt{pr\_mazur}.

\begin{proof}
Use APD2 and the small extension of APD3.  If a nonzero
\(u\in K^\vee\) lay in the kernel of (APD.9), apply its
splitting conclusion to the particular quotient \(C_u\) in
APD4.  Its hypotheses hold exactly: the quotient kills
\(\ker u\), is surjective, and factors the map to \(R\).
The resulting section contradicts APD4.  Therefore
\(K^\vee\hookrightarrow H\), and
\(r=\dim K\le h\).  Choose a basis of \(K\) and lift it
to \(f_1,\ldots,f_r\in I\).  APD3 says these generate
\(I\).  Append \(h-r\) zero elements to obtain the required
family of length \(h\).  This proves both the exact variable
count and the stated bound on generators of the kernel.

No compatibility of the maps \(\mathrm{ob}_p\) for different
small extensions was assumed or needed.  The proof in fact works
over an infinite residue field too; the finite-residue restriction
in the named original statement is certainly preserved.
\end{proof}
% END APD5

% BEGIN APD6
\paragraph{APD6. The strict first-cohomology counting lemma.}
\label{claim:APD6}
Let \(K\) be any field, \(G\) a group with any topology,
\(W,T\) finite-dimensional \(K\)-spaces, \(r\) a
representation of \(G\) on \(W\), and
\(\operatorname{loc}_v:G_v\to G\) any family of group
homomorphisms.  Let \(b_v\subset W\) and
\(\phi:T\to\operatorname{Map}(G,W)\) be linear.  A strict
one-cocycle is a pair \((z,(a_v))\) with \(z\) locally
constant, \(z(st)=z(s)+r(s)z(t)\), and
\[
 z(\operatorname{loc}_v s)-(r(\operatorname{loc}_v s)-1)a_v
                                                        \in b_v.
                                                               \tag{APD.11}
\]
The boundaries are the pairs with \(z(s)=(r(s)-1)X\) and
\(a_v-X\in b_v\).  Quotienting the strict cocycle space by
its intersection with this boundary space defines \(H^1_C\); put
\(H^0_C=W^G\cap\bigcap_v b_v\).
Suppose
\[
 \begin{split}
 &(r(\operatorname{loc}_v s)-1)Y\in b_v\ \text{for all }s
                         \quad\Longrightarrow\quad Y\in b_v,\\
 &\text{for every }(z,a)\text{ as above, there are }D,X
        \text{ with }\phi(D)=z+(r-1)X.
 \end{split}                                                  \tag{APD.12}
\]
Then \(\dim H^1_C\le\dim T\).  Moreover if
\(X_0\ne0\) and \(\phi(D_0)=(r-1)X_0\), then
\[
           \dim H^1_C+1\le\dim T+\dim H^0_C.                 \tag{APD.13}
\]
This is full \texttt{pr\_count}; neither profiniteness nor
finiteness of the field or of the family of local groups is imposed.

\begin{proof}
Let \(C\) be the image of the coboundary map
\(W\to\operatorname{Map}(G,W)\), \(X\mapsto(r-1)X\).
Send a strict cocycle to \([z]\) in
\(\operatorname{Map}(G,W)/C\).  If \(z=(r-1)X\), equation
(APD.11) and the first clause of (APD.12) imply
\(X-a_v\in b_v\) for every \(v\).  Thus the kernel is
exactly the strict boundaries.  The map induced on \(H^1_C\)
is injective, and the second clause of (APD.12) puts its image
inside the image of
\(f:T\to\operatorname{Map}(G,W)/C\).
This proves the first dimension inequality and also proves that
\(H^1_C\) is finite dimensional.

If \(X_0\) is invariant, the first clause of (APD.12) puts it
in every \(b_v\); it is a nonzero element of \(H^0_C\).
The first inequality therefore implies (APD.13).  If it is not
invariant, \(D_0\ne0\), since \(\phi(D_0)\ne0\), but
\(D_0\in\ker f\).  Rank--nullity now gives
\(\dim\operatorname{im}f\le\dim T-1\), again proving
(APD.13).  Local constancy is only used to specify the original
cocycle space; no unstated topology on the coefficient field is used.
\end{proof}
% END APD6

\subsection{Realizing the original local deformation problem}

% BEGIN APD7
\paragraph{APD7. Residue constants in a general complete local ring.}
\label{claim:APD7}
Let \(\mathcal O\to C\) be a local homomorphism, with \(C\)
complete Noetherian local, and suppose its residue field \(k_C\)
is finite as an \(\mathcal O\)-module.  Let \(\kappa\) be
the residue field of \(\mathcal O\).  The subring
\[
 C_0=\{c\in C:\bar c\text{ belongs to the specified image of }
                         \kappa\text{ in }k_C\}
       =\operatorname{im}(\mathcal O)+\mathfrak m_C          \tag{APD.14}
\]
is an \(\mathrm{IsCO}(\mathcal O,C_0)\)-algebra.  The ring
\(C\) is finite as a \(C_0\)-module, and, as subsets of \(C\),
\[
 \mathfrak m_{C_0}=\mathfrak m_C,\qquad
 \mathfrak m_{C_0}^a=\mathfrak m_C^a\quad(a\ge1).
                                                               \tag{APD.15}
\]
For \(a=0\) the corresponding assertion is the intersection
identity \(C_0\cap C=C_0\).  This proves the full original
\texttt{cl\_sub}, without assuming that \(C\) is a DVR or that
its residue constants have a field section in \(C\).

\begin{proof}
Locality of the coefficient map gives an injection
\(\kappa\hookrightarrow k_C\).  The finite-module assumption
therefore says exactly that \([k_C:\kappa]\) is finite.
The inverse image in (APD.14) is a subring.  A member with nonzero
residue has an inverse in \(C\), whose residue is again in
\(\kappa\), so this inverse lies in \(C_0\).  This proves
locality, the asserted maximal ideal, and the prescribed residue
isomorphism.  For \(a\ge1\), a product of \(a\) elements
of \(\mathfrak m_C\), multiplied by an arbitrary coefficient
of \(C\), can absorb that coefficient into its first factor.
All the resulting factors lie in \(\mathfrak m_{C_0}\).
This proves (APD.15) in both directions.

Choose lifts \(b_1=1,b_2,\ldots,b_e\) of a
\(\kappa\)-basis of \(k_C\).  Every element of \(C\)
differs from a combination of the \(b_i\), with coefficients
lifted from \(\mathcal O\), by an element of
\(\mathfrak m_C\subset C_0\).  Thus the \(b_i\) generate
\(C\) over \(C_0\).  If \(x_j\) generate \(\mathfrak m_C\)
as a \(C\)-ideal, the finite family \(b_ix_j\) generates it
as a \(C_0\)-module.

Here is also a direct proof of Noetherianity of \(C_0\).
For an ideal \(I\subset C_0\), choose finitely many elements
of \(I\) generating \(CI\) as a \(C\)-ideal.  This makes
\(CI\) finite over \(C_0\).  The product
\(\mathfrak m_{C_0}CI\) is a finite \(C_0\)-module and is
contained in \(I\): indeed \(\mathfrak m_{C_0}C\subset C_0\).
The quotient \(CI/\mathfrak m_{C_0}CI\) is a finite-dimensional
\(\kappa\)-space.  Its subspace
\(I/\mathfrak m_{C_0}CI\) has a finite basis, and lifts of
that basis, together with generators of
\(\mathfrak m_{C_0}CI\), generate \(I\).  Hence every ideal
is finite.  Finally \(C_0\) is closed in \(C\), being the
inverse image of a subset of the discrete residue field.
Equation (APD.15) identifies its adic topology with the induced
one; completeness and separation follow.  All its coefficient
structures are restrictions of the original map.
\end{proof}
% END APD7

% BEGIN APD8
\paragraph{APD8. Broad residual points have one triangular direction.}
\label{claim:APD8}
Let \(K\) be any field, \(f:k\to K\), and let
\(r,\psi\) be an original normalized rank-two or rank-three
representation satisfying \(\mathrm{IsBar}(f)\).  This condition
prescribes the diagonal entries and ordinary characters; it does
not prescribe the off-diagonal entries.  Nevertheless either all
matrices \(r(g)\) are upper triangular or all are lower triangular.
Consequently a weight-zero monomial in their entries, containing
an off-diagonal factor, is zero.

\begin{proof}
The diagonal entries are characters
\(f(\Psi(g)^{w_i})\), where
\(w=(1,-1)\) or \((1,0,-1)\).  Comparing the \((0,0)\)
entry of \(r(gh)=r(g)r(h)\) in rank two gives
\(r(g)_{01}r(h)_{10}=0\) for all \(g,h\).  Over a field,
either the entire first function is zero or the entire second one
is zero.  These are the two alternatives.

In rank three use the normalized marked difference
\[
 Z=r(g_0)-r(g_0^*)=\operatorname{diag}(q,0,-q),
 \qquad q=f(\Psi(g_0)-\Psi(g_0)^{-1}).                       \tag{APD.16}
\]
The frame gives \(q\ne0\), and its odd characteristic gives
\(q\ne-q\).  For every \(i\), the diagonal character identity
applied to \(gh\) and to
\(gg_0h,gg_0^*h\) gives
\[
 \sum_{j\ne i}r(g)_{ij}r(h)_{ji}=0,
 \qquad
 \sum_{j\ne i}Z_{jj}r(g)_{ij}r(h)_{ji}=0.                 \tag{APD.17}
\]
There are only two summands, with distinct coefficients in the
second equation.  Subtraction after multiplying the first equation
therefore makes each summand zero.  In particular the pairs of
entry functions \((01,10)\), \((02,20)\), and \((12,21)\)
cannot both be nonzero.

Polarization gives \(r(g^*)_{01}=r(g)_{12}\) and
\(r(g^*)_{10}=r(g)_{21}\), so the two adjacent upper functions
vanish together, as do the two adjacent lower functions.  If
\(01\) vanishes, so does \(12\).  If \(02\) also vanishes
the representation is lower triangular.  Otherwise \(20\)
vanishes; the \((0,1)\) entry of the product identity, with
\(01=12=20=0\), gives
\(r(g)_{02}r(h)_{21}=0\).  It follows that \(21=10=0\),
so the representation is upper triangular.  The branch in which
\(10\) vanishes is the same argument with upper and lower
indices exchanged.  Thus one direction works for all group
elements at once.

The weights strictly decrease with the index.  In an upper
triangular monomial all factors that can be nonzero have
nonnegative weight, and every off-diagonal such factor has
positive weight.  A weight-zero monomial containing an
off-diagonal factor therefore has a zero factor.  The lower
triangular case reverses the signs.  This proof works over every
specified target field and divides only by nonzero differences
and by the integer \(2\); the allowed Host case \(\ell=3\)
is unchanged.
\end{proof}
% END APD8

% BEGIN APD9
\paragraph{APD9. The graded Laurent vertex is an actual separated algebra.}
\label{claim:APD9}
Suppose \(C\) is an \(\mathrm{IsCO}(O,C)\)-algebra,
\(C\hookrightarrow K\) is an injective map of commutative
rings with the specified \(O\)-scalar tower, and
\(\sigma=(r,\psi):\mathrm{NRep}(F,n,K)\) satisfies the original
\(\mathrm{Cone}(O,C,\sigma)\).  Explicitly, its diagonal
entries and full ordinary character values lie in \(C\) with
the prescribed residues; every weight-zero entry monomial with
an off-diagonal factor lies in \(\mathfrak m_C\); each entry
span \(M_{ij}=\langle r(g)_{ij}:g\in G_0\rangle_C\) is
finite; the entries are locally constant modulo
\(\mathfrak m_C^aM_{ij}\); and the character functions in
\(C\) are adically continuous.

Let \(P_j\subset K\) be the \(C\)-span of weight-\(j\)
entry monomials and define
\[
 D=\bigoplus_{j\in\mathbb Z}P_jt^j\subset K[t,t^{-1}],
 \qquad
 \mathfrak n=\{d\in D:[t^0]d\in\mathfrak m_C\}.            \tag{APD.18}
\]
Then \(D\) is Noetherian, \(\mathfrak n\) is a maximal ideal,
\(D/\mathfrak n=\kappa=k\), and
\(\bigcap_a\mathfrak n^a=0\).  There is a normalized
representation \(\sigma_D\) over \(D\) with
\[
 r_D(g)_{ij}=t^{w_i-w_j}r(g)_{ij},\qquad
 \psi_{D,v,i}(s)=\psi_{v,i}(s).                            \tag{APD.19}
\]
It is \(\mathfrak n\)-adically continuous and has the original
diagonal residual lift.

\begin{proof}
The products of the \(P_j\)'s satisfy \(P_iP_j\subset P_{i+j}\).
The degree-zero part is exactly \(C\): a diagonal-only monomial
lies in \(C\), and any other degree-zero monomial lies in
\(\mathfrak m_C\) by the cone condition.  For \(j\ne0\),
a product from \(P_jP_{-j}\) has an off-diagonal factor, so
lies in \(\mathfrak m_C\).  Thus taking constant coefficient
modulo \(\mathfrak m_C\) is a ring homomorphism
\(D\to k\), with kernel exactly \(\mathfrak n\).

For each pair \((i,j)\), choose finitely many generators of
\(M_{ij}\) and multiply them by \(t^{w_i-w_j}\).  This finite
family generates \(D\) as a \(C\)-algebra: every weighted
entry is in its span, every weighted monomial is a product of
these entries, and a Laurent polynomial has finite support.
Hilbert's basis theorem makes \(D\) Noetherian.

We spell out separation, since the vertex ideal need not be in
the Jacobson radical of the whole Laurent subalgebra.  It is a
homogeneous ideal, so its powers and their intersection
\(J=\bigcap_a\mathfrak n^a\) are homogeneous.  Artin--Rees
for the finite ideal \(J\subset D\) gives \(J=\mathfrak nJ\).
The determinant trick applied to finite generators of \(J\)
then gives \(b\in\mathfrak n\) with \((1-b)J=0\): writing
each generator as an \(\mathfrak n\)-linear combination of
the generators, the determinant of the corresponding
identity-minus-coefficient matrix is \(1-b\) and annihilates
them.  If \(at^j\in J\) is homogeneous, the degree-\(j\)
part of \(b(at^j)=at^j\) says \(ca=a\), where
\(c=[t^0]b\in\mathfrak m_C\).  The image of \(1-c\) is
a unit in \(K\), hence \(a=0\).  Every homogeneous part of
every element of \(J\) is zero, so \(J=0\).

Formula (APD.19) is conjugation by
\(\operatorname{diag}(t^{w_i})\) in \(K[t,t^{-1}]\).
Its inverse matrices come from \(g^{-1}\), and each ordinary
character has its inverse at \(s^{-1}\) already in \(C\).
The normalized marking, polarization, and the ordered ordinary
product identities therefore hold in \(D\); the relation
\(w_{i^{\mathrm{rev}}}=-w_i\) checks polarization explicitly.
An entry difference in \(\mathfrak m_C^aM_{ij}\), after
weighting, is a sum of elements of \(\mathfrak n^a\) times
weighted generators of \(M_{ij}\).  The cone continuity
condition consequently gives continuity modulo every
\(\mathfrak n^a\).  Characters use the same argument in
degree zero.  Every off-diagonal weight is nonzero, so its
weighted entry lies in \(\mathfrak n\); the diagonal and
character residues are exactly the specified ones.
\end{proof}
% END APD9

% BEGIN APD10
\paragraph{APD10. Cone points and broad local points are realized.}
\label{claim:APD10}
At the hypotheses of APD9 there are actual \(O\)-algebra maps
\[
 \xi:\mathcal A\longrightarrow K,
 \qquad\beta:\mathcal B\longrightarrow C,
 \qquad \xi_*u=\sigma,\qquad
 \xi(\mathrm{inc}(b))=\beta(b)\text{ in }K.                \tag{APD.20}
\]
This proves full original \texttt{alg}.  Consequently, for every
complete Noetherian local \(O\)-algebra \(C\) satisfying the
original \(\mathrm{IsCL}(O,C)\), every compatible residue map
\(f:k\to k_C\), and every adically continuous normalized
\(\sigma\) with \(\bar\sigma.\mathrm{IsBar}(f)\), there is
\[
 \xi:\mathcal A\longrightarrow C,
 \qquad \xi_*u=\sigma,
 \qquad \xi(\mathfrak m_{\mathcal B}\mathcal A)
                                  \subset\mathfrak m_C. \tag{APD.21}
\]
Thus the full original \texttt{point\_loc} is supplied, with its
broader residual condition intact.

\begin{proof}
Use the vertex algebra of APD9.  Each
\(D_a=D/\mathfrak n^{a+1}\) is Noetherian local, has nilpotent
maximal ideal, and has the prescribed residue field \(k\).
It is therefore a complete coefficient algebra over \(O\).
The universal property included in \(S\) gives a unique map
\(f_a:S.R\to D_a\) carrying the universal representation to
\(\sigma_D\bmod\mathfrak n^{a+1}\).  Uniqueness makes these
maps compatible.  They kill \(\mathfrak m_{S.R}^{a+1}\), since
coefficient-algebra maps are local.

Let \(\widehat D=\varprojlim D_a\).  The resulting map
\(f:S.R\to\widehat D\) sends every matrix entry to its actual
element of \(D\).  It sends every weight-zero monomial and
ordinary character value to \(C\).  We must also handle their
closure defining \(\mathcal B\).  The complete Noetherian ring
\(C\) has finite residue field, hence each
\(C/\mathfrak m_C^a\) is finite and \(C\) is compact in its
adic topology.  Its map to the Hausdorff ring \(\widehat D\)
is continuous and injective, the latter by APD9.  Its image is
therefore closed.  Approximating an element of \(\mathcal B\)
by polynomials in its defining generators, and using that \(f_a\)
kills \(\mathfrak m_{S.R}^{a+1}\), shows that its image lies
in this closed copy of \(C\).  This is a concrete coefficient
limit argument; it does not presume that \(D\) is complete.

Since \(\mathcal A\) is generated by \(\mathcal B\) and
the actual matrix entries, \(f(\mathcal A)\subset D\).
Separation of \(D\) makes this a unique algebra map into
\(D\), and its restriction to \(\mathcal B\) gives \(\beta\).
Evaluation \(t=1\) is a ring homomorphism on finite Laurent
polynomials and sends (APD.19) to \(\sigma\).  Composing
gives (APD.20).  In particular no evaluation of an unrestricted
completed Laurent series has been used.

For (APD.21), take \(C_0\subset C\) from APD7.  The
\(\mathrm{IsCL}\) hypothesis supplies precisely its finite
residue-extension premise.  The residual diagonal and character
conditions put all those values, and their specified residues,
in \(C_0\).  APD8 puts each weight-zero monomial containing an
off-diagonal entry in \(\mathfrak m_C=\mathfrak m_{C_0}\).
Each entry span over \(C_0\) is finite because \(C_0\) is
Noetherian and \(C\) is finite over it.

For clarity, its required entry-span continuity does not follow
just by renaming the topology.  If
\(M\subset C\) is one of the finite entry spans, Artin--Rees
(ALC6) gives, for each \(a\), an integer \(b\) with
\[
 M\cap\mathfrak m_{C_0}^{b}C
                         \subset\mathfrak m_{C_0}^{a}M.
                                                               \tag{APD.22}
\]
An entry difference in \(\mathfrak m_C^b\) belongs to the
left side, by (APD.15), so adic continuity in \(C\) gives the
cone's exact quotient-module continuity.  Character continuity
in \(C_0\) follows directly from (APD.15), including their
inverse values.  All five original cone clauses are now proved.
Apply (APD.20) with \(C_0\hookrightarrow C\).
The map \(\beta:\mathcal B\to C_0\) preserves residue
maps, since RTC7 makes \(\mathcal B\) an
\(\mathrm{IsCO}(O,-)\)-algebra.  It sends
\(\mathfrak m_{\mathcal B}\) into \(\mathfrak m_{C_0}\).
This proves the last assertion of (APD.21) on the original
inclusion \(\mathcal B\to\mathcal A\).
\end{proof}
% END APD10

% BEGIN APD11
\paragraph{APD11. Entries determine the map on the actual entry algebra.}
\label{claim:APD11}
Let \(C\) be any Noetherian local \(O\)-algebra; it need not
be complete.  If two maps \(\xi,\xi':\mathcal A\to C\)
carry \(u\) to the same normalized representation and both send
\(\mathfrak m_{\mathcal B}\mathcal A\) into \(\mathfrak m_C\),
then \(\xi=\xi'\).  This is the full original \texttt{ext\_A}.

\begin{proof}
The maps agree on matrix entries, ordinary character values and
their inverses, and hence on the algebra generated by weight-zero
entry monomials and character values.  Fix \(b\in\mathcal B\)
and \(a\ge1\).  RTC9's fixed-weight comparison at weight zero
gives \(b_0\) in that generating algebra with
\(b-b_0\in\mathfrak m_{\mathcal B}^{a}\mathcal B\): choose
the required ambient \(\mathfrak m_{S.R}\)-adic approximation
and then apply the comparison.  Both maps send this difference
into \(\mathfrak m_C^a\).  Their values on \(b_0\) agree,
so \(\xi(b)-\xi'(b)\in\mathfrak m_C^a\) for every \(a\).
Krull intersection in the Noetherian local ring \(C\) makes
the difference zero.  They thus agree on \(\mathcal B\) and
on all matrix-entry generators of \(\mathcal A\), proving
the assertion.  This uses the actual fixed-weight topology; no
continuity of an arbitrary map out of \(S.R\) is assumed.
\end{proof}
% END APD11

% BEGIN APD12
\paragraph{APD12. The actual local representation is universal.}
\label{claim:APD12}
For the supplied \(L\) in (APD.1), let \(C\) be every
\(\mathrm{IsCO}(O',C)\)-algebra and let \(\sigma\) be every
continuous normalized representation with
\[
 \operatorname{res}_{C,*}\sigma
                     =\operatorname{res}_{R,*}\rho
                    \quad\text{over }k'=\operatorname{Res}(O').
                                                               \tag{APD.23}
\]
There is a unique \(O'\)-algebra map \(g:R\to C\) with
\(g_*\rho=\sigma\).  Together with ALC17's continuity,
this is the full original \texttt{def\_iso} and the literal
\(\rho.\mathrm{Univ}(O')\) required by \texttt{pres}.

\begin{proof}
The map \(O\to O'\) is local: the original unramified
identity \(\mathfrak m_OO'=\mathfrak m_{O'}\) proves this
directly.  Its composite to \(C\) is therefore local.
The ring \(C\) is complete Noetherian, and its residue field
is \(k'\), finite over \(O\) because \(O'\) is a finite
\(O\)-module.  Thus it satisfies the precise
\(\mathrm{IsCL}(O,C)\) input of APD10.

The compatible maps \(k\to k'\to\operatorname{Res}(C)\)
and \(k\to k'\to\operatorname{Res}(R)\) are those induced
by the specified scalar towers.  ALC17 gives the residual
\(\mathrm{IsBar}\) condition for \(\rho\).  Composing
(APD.23) with the second residue identification gives the same
condition for \(\sigma\).  APD10 produces
\(\xi:\mathcal A\to C\), with \(\xi_*u=\sigma\), sending
\(x_0=\mathfrak m_{\mathcal B}\mathcal A\) into
\(\mathfrak m_C\).

Compare \(\operatorname{res}_C\xi\) and
\(\operatorname{res}_R\iota\) as maps \(\mathcal A\to k'\).
Their normalized representations agree by (APD.23).  Both kill
\(x_0\): for the latter this follows from
\(x_0\subset x=\iota^{-1}(\mathfrak m_R)\), proved in
ALC11.  APD11 therefore says these two residue maps are equal.
The original universal field of the supplied \(L\) now gives
a unique \(O'\)-map \(g:R\to C\) extending \(\xi\).
It sends \(\rho\) to \(\sigma\).

If \(g'\) is any other map sending the representation correctly,
the maps \(g'\iota\) and \(\xi\) agree on the representation
and both send \(x_0\) into \(\mathfrak m_C\), because
\(g'\) is a coefficient-algebra map.  APD11 makes them equal,
and the uniqueness field of \(L\) gives \(g'=g\).
The same \(O',R,\iota,x\) was fixed before \(C\) and
\(\sigma\) were chosen.  No diagonal residual hypothesis was
added to (APD.23).
\end{proof}
% END APD12

\subsection{Continuous boxes and square-zero corrections}

\paragraph{Scope and notation for the deformation lemmas.}
In APD13--APD20 the coefficient ring \(\mathcal O\) is an
arbitrary local ring, \(\kappa=\operatorname{Res}(\mathcal O)\),
and the rank \(n\) is arbitrary unless explicitly restricted.
The frame is any original \(F\) over a field of characteristic
\(\ell\), with \(F.\mathrm{Ok}(C_0)\).  This broader scope is
the scope of the original matrix lemmas.  It is not limited to
the arithmetic \(O',R\) of (APD.1).
Put
\[
 \theta(A)=J_n(A^{-1})^tJ_n,\qquad
 \theta_{\mathrm{Lie}}(X)=-J_nX^tJ_n,
 \qquad \mathrm{GG}_n(C)=\mathrm{GL}_n(C)\rtimes\{\pm1\},
                                                               \tag{APD.24}
\]
where the negative sign acts by \(\theta\).  A box
\(\mathfrak B\) consists of a homomorphism
\(b:\Gamma\to\mathrm{GG}_n(C)\), with sign \(F.\delta\),
and matrices \(g_v\) making the restriction to each \(G_v\)
upper triangular.  It is \emph{of} a normalized representation
\(\rho\) if its matrix at \(c\) is \(1\), its restriction
to \(G_0\) is \(\rho.r\), and the diagonal in the \(g_v\)
basis is the specified ordered character family \(\rho.\psi\).
This is exactly \(\mathrm{Box.Of}\), not a change of deformation
problem.  On \(W=M_n(\kappa)\) its residual adjoint action is
\[
 \overline{\operatorname{Ad}}(s)X=
 \begin{cases}
  \bar b(s)X\bar b(s)^{-1},&\delta(s)=1,\\
  \bar b(s)\theta_{\mathrm{Lie}}(X)\bar b(s)^{-1},&\delta(s)=-1,
 \end{cases}
 \quad
 \mathfrak b_v=\{X:\bar g_vX\bar g_v^{-1}
                                    \text{ is upper triangular}\}.
                                                               \tag{APD.25}
\]
Here \(\bar b(s)\) denotes the matrix part.  Bars on boxes
and representations use the specified residue map of the
coefficient algebra.

% BEGIN APD13
\paragraph{APD13. Boxes are constructed and recovered with their characters.}
\label{claim:APD13}
For every commutative ring \(C\), ideal \(J\subset C\), and
\(J\)-adically continuous normalized representation \(\rho\)
with an ordered ordinary flag at every \(v\), there is a box
\(\mathfrak B\) of \(\rho\), with \(J\)-adically continuous
matrix entries on \(\Gamma\).  Conversely a continuous box
with matrix \(1\) at \(c\) and diagonal marked difference
comes from a continuous normalized representation, with its full
ordered ordinary characters.  These are full original
\texttt{pr\_box} and \texttt{pr\_ofbox}.

\begin{proof}
Since \(c^2=1\) and \(\delta(c)=-1\), every \(s\in\Gamma\)
has a unique expression \(h\) or \(hc\), with \(h\in G_0\).
Define the matrix part of \(b(s)\) to be \(\rho.r(h)\) and
its sign to be \(\delta(s)\).  Polarization gives
\(\rho.r(chc^{-1})=\theta(\rho.r(h))\); this proves the
semidirect-product homomorphism law in the four possible sign
cases.  Its value at \(c\) has matrix \(1\).  Choose the
given ordinary bases \(g_v\); they give the required local
upper triangularity and ordered diagonal entries.

The map \(s\mapsto h\) is continuous: it is \(s\) on the
open subgroup \(G_0\) and \(sc\) on its open other coset.
Thus continuity of \(\rho\) proves the stated entry continuity
of \(b\), at every power of \(J\).

Conversely restrict the matrix part of the box to \(G_0\),
where the sign is positive.  Conjugate locally by \(g_v\),
and take its diagonal entries as characters.  Diagonal entries
of a product of upper triangular matrices multiply, and those
of the inverse are inverse units, so these really are characters.
An upper triangular matrix with its \(i\)-th diagonal entry
subtracted takes the \((i+1)\)-st standard flag space into the
\(i\)-th.  Starting with the rightmost factor proves the
ordered ordinary-product identity for every tuple of local
group elements.  The normalization at \(c\) gives
\(r(h^*)=J_nr(h)^tJ_n\), and the given marked difference
gives the remaining normalization.  Restriction and finite matrix
sums prove continuity of both matrix entries and character values.
Reduction and arbitrary coefficient maps preserve all of these
constructions.
\end{proof}
% END APD13

% BEGIN APD14
\paragraph{APD14. Continuous set-theoretic lifts across every small extension.}
\label{claim:APD14}
Let \(p:C'\twoheadrightarrow C\) be a surjection of
\(\mathrm{IsCO}(\mathcal O,-)\)-algebras, with
\(K=\ker p\) and \(\mathfrak m_{C'}K=0\).  Every
adically continuous function from any topological space into
\(C\) has an adically continuous lift to \(C'\).
Consequently every continuous box \(\mathfrak B\) over \(C\)
has the original \(\mathrm{SetLift}(p,\mathfrak B)\): continuous
matrix lifts \(r_s\) of its matrices, lifts \(g'_v\) of its
ordinary bases, and continuous upper triangular lifts \(t_{v,s}\)
of \(g_vb(\operatorname{loc}_v s)g_v^{-1}\).
This proves full \texttt{pr\_sec} and \texttt{pr\_setlift}.

\begin{proof}
The map preserves residues, so
\(p^{-1}(\mathfrak m_C)=\mathfrak m_{C'}\) and
\(p(\mathfrak m_{C'}^a)=\mathfrak m_C^a\) for every \(a\).
Artin--Rees applied to \(K\subset C'\), together with
\(\mathfrak m_{C'}K=0\), gives an integer \(a_0\ge1\)
with \(K\cap\mathfrak m_{C'}^{a_0}=0\).  Thus
\(p:\mathfrak m_{C'}^{a_0}\to\mathfrak m_C^{a_0}\) is an
additive bijection.  Choose a representative in \(C'\) for
each class of \(C/\mathfrak m_C^{a_0}\).  Lift a general
\(z\in C\) by its chosen representative plus the unique
element of \(\mathfrak m_{C'}^{a_0}\) lifting the remaining
difference.  This defines a section \(s:C\to C'\).

If \(z-z'\in\mathfrak m_C^a\) with \(a\ge a_0\), their
chosen representatives coincide.  The difference \(s(z)-s(z')\)
belongs to \(\mathfrak m_{C'}^{a_0}\) and lifts \(z-z'\).
There is also a lift in \(\mathfrak m_{C'}^a\); uniqueness
in \(\mathfrak m_{C'}^{a_0}\) makes these equal.  Hence
\(s(z)-s(z')\in\mathfrak m_{C'}^a\).  Applying this to a
continuous function at level \(\max(a,a_0)\) proves its lifted
continuity at level \(a\).  No compactness or profiniteness of
its domain is required.

Apply these lifts entry by entry to an invertible matrix-valued
function.  Its lifted determinant is a unit, since its image is
a unit under the local map \(p\), so the lifted matrix is
invertible.  If an entry function is identically zero, choose
its lift identically zero.  This preserves upper triangularity
for the local matrices.  Lift each fixed \(g_v\) in the same
way.  The local conjugated matrices are continuous because the
local group maps and finite matrix operations are continuous
modulo every power.  These choices give every field of
\(\mathrm{SetLift}\).  Neither the lifted \(r_s\) nor the
lifted \(t_{v,s}\) is asserted to be a homomorphism.
\end{proof}
% END APD14

% BEGIN APD15
\paragraph{APD15. The two explicit defects are continuous cone cocycles.}
\label{claim:APD15}
Let \(p:C'\to C\) be any map of coefficient algebras with
\(\mathfrak m_{C'}K=0\), and suppose a set lift
\(D=(r_s,g'_v,t_{v,s})\) of a box is supplied.  Define
\[
 \begin{split}
 \Omega(s,t)&=r_s\,\theta^{\delta(s)}(r_t)r_{st}^{-1}-1,\\
 A_v(s)&=r_{\operatorname{loc}_v s}
               ((g'_v)^{-1}t_{v,s}g'_v)^{-1}-1,
 \end{split}                                                  \tag{APD.26}
\]
where \(\theta^{+1}\) means the identity and
\(\theta^{-1}\) means \(\theta\).  All their entries belong
to \(K\), and both functions are locally constant.  For every
\(\mathcal O\)-linear \(u:K\to\kappa\), the entrywise
images \((\omega_u,a_{u,v})\) satisfy
\[
 \begin{split}
 \overline{\operatorname{Ad}}(s)\omega_u(t,r)+\omega_u(s,tr)
                 &=\omega_u(st,r)+\omega_u(s,t),\\
 \omega_u(\operatorname{loc}_v s,\operatorname{loc}_v t)
      -d a_{u,v}(s,t)&\in\mathfrak b_v,\\
 d\beta(s,t)&=\overline{\operatorname{Ad}}(s)\beta(t)
                                  -\beta(st)+\beta(s).
 \end{split}                                                  \tag{APD.27}
\]
The assignment \(u\mapsto(\omega_u,a_u)\) is linear.
This proves full \texttt{pr\_omk},
\texttt{borelProjection\_\allowbreak isTwoCocycle}, and \texttt{pr\_cocv}.

\begin{proof}
Mapping (APD.26) by \(p\) gives zero, by the homomorphism law
of the box and the three defining map identities of a set lift.
Also \(K\subset\mathfrak m_{C'}\), so \(K^2=0\).
Artin--Rees as in APD14 gives \(a_0\) with
\(K\cap\mathfrak m_{C'}^{a_0}=0\).  The functions in (APD.26)
are continuous modulo every power, since inversion in the reduced
matrix group is a function of the reduced invertible matrix and
\(\delta\) is locally constant.  Equality modulo
\(\mathfrak m_{C'}^{a_0}\) for two values in \(M_n(K)\)
is actual equality.  Their functions are therefore locally
constant, without any assumption that \(\kappa\) is finite.

The action of a coefficient \(c\in C'\) on \(K\) is its
residue action: subtract a residue lift from \(\mathcal O\)
and use \(\mathfrak m_{C'}K=0\).  Thus \(u(cy)=\bar c\,u(y)\).
For matrices \(X,Y\) with entries in \(K\),
\((1+X)(1+Y)=1+X+Y\) and \((1+X)^{-1}=1-X\).
Conjugation by a lifted signed matrix acts on these entries
through (APD.25), since
\(\theta(1+X)=1+\theta_{\mathrm{Lie}}(X)\).

Put \(x_s=(r_s,\delta(s))\).  Associativity gives the exact
identity
\[
 (x_sx_tx_{st}^{-1})(x_{st}x_rx_{str}^{-1})
 =x_s(x_tx_rx_{tr}^{-1})x_s^{-1}(x_sx_{tr}x_{str}^{-1}).
\]
Substitution of \(1+\Omega\), followed by \(K^2=0\), gives
the first equation of (APD.27), after applying \(u\).
For the local equation put
\(T_v(s,t)=t_{v,s}t_{v,t}t_{v,st}^{-1}-1\).
It has entries in \(K\) and is upper triangular.  Substitution
of the definitions and cancellation of the invertible factors gives
\[
 \Omega(\operatorname{loc}_v s,\operatorname{loc}_v t)
 -\bigl(r_{\operatorname{loc}_v s}A_v(t)
             r_{\operatorname{loc}_v s}^{-1}
                 -A_v(st)+A_v(s)\bigr)
                 =(g'_v)^{-1}T_v(s,t)g'_v.                \tag{APD.28}
\]
One can check this without commuting matrices: multiply
\((1+\Omega)(1+A_v(st))\), replace the lifted local matrices
by \((1+A_v(s))(g'_v)^{-1}t_{v,s}g'_v\), and cancel.
All products of two kernel matrices disappear.  Applying \(u\)
to (APD.28) makes its right side an element of
\(\mathfrak b_v\), proving the second equation.  Finally all
applications of \(u\) are to actual kernel entries; their
linearity proves linearity of the entire assignment.
\end{proof}
% END APD15

% BEGIN APD16
\paragraph{APD16. A vanishing two-cocycle gives the actual corrected box.}
\label{claim:APD16}
Let \(\mu:\kappa\to C'\) be an additive map satisfying
\[
 \mu(\bar c a)=c\mu(a),\qquad
 \mu(\kappa)\subset\mathfrak m_{C'},\qquad p\mu=0.
                                                               \tag{APD.29}
\]
Here \(p:C'\to C\) is any map between the original
\(\mathrm{IsCO}\)-algebras, and the set lift is any original
\(\mathrm{SetLift}(p,\mathfrak B)\).  We do not require
\(\mathfrak m_{C'}\ker p=0\), nor surjectivity of \(p\).
Only the image of \(\mu\) must be square zero, as follows
from (APD.29).  Suppose \(\beta:\Gamma\to W\)
is locally constant, \(Y_v\in W\), and functions \(a_v\)
satisfy
\[
 \Omega=\mu(d\beta),\qquad A_v=\mu(a_v),\qquad
 a_v(s)-\beta(\operatorname{loc}_v s)
       +(\overline{\operatorname{Ad}}(\operatorname{loc}_v s)-1)Y_v
                                                        \in\mathfrak b_v.
                                                               \tag{APD.30}
\]
Entrywise application of \(\mu\) is understood.  The formulas
\[
 b'(s)=\bigl(1-\mu(\beta(s))\bigr)(r_s,\delta(s)),
 \qquad g'_v{}^{\mathrm{new}}=g'_v(1-\mu(Y_v))             \tag{APD.31}
\]
define a continuous box over \(C'\) whose image is the given box
over \(C\).  These are the actual corrections in full original
\texttt{pr\_sqz} and \texttt{pr\_correct}.

\begin{proof}
Equation (APD.29) implies \(\mu(a)\mu(b)=0\), since the
residue of \(\mu(a)\) is zero.  Consequently
\(E(X)=1+\mu(X)\) is invertible,
\(E(X)E(Y)=E(X+Y)\), and \(E(X)^{-1}=E(-X)\).
For any signed invertible matrix \(z\), (APD.29) gives
\[
 zE(X)z^{-1}=E(\operatorname{Ad}(\bar z)X).                \tag{APD.32}
\]
This follows from matrix multiplication for positive sign; for
negative sign use \(\theta(E(X))=E(-J_nX^tJ_n)\).
The same coefficient calculation gives, for every matrix \(A\)
over \(C'\),
\(A\mu(X)=\mu(\bar A X)\) and
\(\mu(X)A=\mu(X\bar A)\), by expanding each matrix entry.
These identities and (APD.32) use only the first two conditions
of (APD.29).  They therefore give all outputs of
\texttt{pr\_sqz} for any \(\mathrm{IsCO}(\mathcal O,C')\),
any rank, and any original \(\mathrm{IsSq}\) map, without a
frame, a map \(p\), or its kernel conditions.

Multiply the two expressions for \(b'(s)\) and \(b'(t)\)
using (APD.32).  Their correction exponent is
\(-\beta(s)-\overline{\operatorname{Ad}}(s)\beta(t)
 +d\beta(s,t)=-\beta(st)\).  Thus \(b'\) is a group
homomorphism, with the unchanged sign character.  For a local
element, conjugating by the new \(g'_v\) gives
\[
 E\bigl(\bar g_v[
     a_v(s)-\beta(\operatorname{loc}_v s)
       +(\overline{\operatorname{Ad}}(\operatorname{loc}_v s)-1)Y_v
                ]\bar g_v^{-1}\bigr)t_{v,s}.              \tag{APD.33}
\]
The bracket condition in (APD.30) says its exponential is upper
triangular.  The product is therefore upper triangular as well.
Both the representation and bases reduce to the given ones because
\(p\mu=0\).  The new matrix entries are finite sums of
continuous entries times values of the locally constant \(\beta\).
They are continuous modulo every maximal-ideal power.  This
verifies every box field and the exact formula (APD.31), including
its signs.
\end{proof}
% END APD16

% BEGIN APD17
\paragraph{APD17. Both normalization equations are solved explicitly.}
\label{claim:APD17}
Let \(p:C'\to C\) be any map of coefficient algebras with
\(\mathfrak m_{C'}\ker p=0\).  Let \(\rho\) be normalized
over \(C\), and suppose its residue satisfies
\(\rho.\mathrm{Good}\): \(2\ne0\), the diagonal entries
of its marked difference are pairwise distinct, and each local
character list has a separating group element.  If a continuous
box \(\mathfrak B'\) maps to a box of \(\rho\), then there
are \(h\in\mathrm{GL}_n(C')\), with \(p(h)=1\), and a
continuous normalized \(\rho'\) over \(C'\), with
\(p_*\rho'=\rho\), such that the conjugated box has matrix
\(1\) at \(c\) and restricts to \(\rho'\) on \(G_0\).
This proves full \texttt{pr\_norm1}, \texttt{pr\_norm2}, and
\texttt{pr\_norm}; no surjectivity of \(p\) is required.
More precisely, the two elementary assertions below also allow
an arbitrary target \(\mathcal O\)-algebra \(C\), with no
\(\mathrm{IsCO}\) structure: assume explicitly
\(\ker p\subset\mathfrak m_{C'}\).  The first takes any
signed involution of negative sign with matrix reducing to
\(1\) and \(2\ne0\) in \(\kappa\).  The second takes
any \(Z\) fixed by \(\theta_{\mathrm{Lie}}\), with all
off-diagonal entries in \(\ker p\) and pairwise distinct
residual diagonal entries.  Both retain
\(\mathfrak m_{C'}\ker p=0\).  These are the literal broader
premises of \texttt{pr\_norm1} and \texttt{pr\_norm2}.

\begin{proof}
Write \(K=\ker p\).  It lies in \(\mathfrak m_{C'}\), so
\(K^2=0\); moreover \(2\) is a unit in \(C'\).  If the
matrix at \(c\) is \(1+X\), then \(X\in M_n(K)\).
The involution identity at \(c\) gives
\(J_nX^tJ_n=X\).  Put \(h_1=1-X/2\).  Its inverse is
\(1+X/2\), and signed conjugation gives
\[
 h_1(1+X)\theta(h_1^{-1})
                 =(1-X/2)(1+X)(1-X/2)=1.                 \tag{APD.34}
\]
Thus the first normalization is solved by an explicit kernel
matrix, with no lifting of an orthogonal group assumed.

In the resulting box put \(Z=b(g_0)-b(g_0^*)\).  Polarization
now gives \(\theta_{\mathrm{Lie}}(Z)=Z\).  Its off-diagonal
entries lie in \(K\), and the differences \(Z_{ii}-Z_{jj}\)
are units for \(i\ne j\), by the residual separation condition.
Define
\[
 A_{ii}=0,\qquad A_{ij}=Z_{ij}(Z_{ii}-Z_{jj})^{-1}\ (i\ne j),
 \qquad h_2=1+A.                                         \tag{APD.35}
\]
Every entry of \(A\) belongs to \(K\), so \(A^2=0\).
The relation \(Z_{j^{\mathrm{rev}},i^{\mathrm{rev}}}=-Z_{ij}\)
gives \(A_{j^{\mathrm{rev}},i^{\mathrm{rev}}}=-A_{ij}\);
therefore \(\theta(h_2)=h_2\).  This conjugation preserves
the matrix \(1\) at \(c\).  For \(i\ne j\), expansion gives
\[
 ((1+A)Z(1-A))_{ij}
  =Z_{ij}+A_{ij}Z_{jj}-Z_{ii}A_{ij}=0.                    \tag{APD.36}
\]
The omitted terms each contain two entries of \(K\) and are
zero; this explains why the formula works over an arbitrary
coefficient algebra, not just a field.  It proves the second
normalization equation.

Conjugate each local basis simultaneously by replacing it with
\(g'_v h_1^{-1}h_2^{-1}\).  Upper triangularity and its
ordered diagonal characters are thereby preserved.  APD13 now
recovers a normalized \(\rho'\).  All conjugating matrices
reduce to \(1\), so after applying \(p\) this box is the
original one; its restricted representation and diagonal
characters are precisely \(\rho\).  Constant conjugations
and the locally constant sign character preserve all adic
continuity assertions.  The two elementary calculations also
prove \texttt{pr\_norm1} and \texttt{pr\_norm2} at their
broader literal premises: they only need the stated kernel,
involution and diagonal conditions, respectively.
\end{proof}
% END APD17

% BEGIN APD18
\paragraph{APD18. The continuous cone is the original obstruction space.}
\label{claim:APD18}
Let \(R\) be an \(\mathrm{IsCO}(\mathcal O,R)\)-algebra,
let \(\rho.\mathrm{Univ}(\mathcal O)\) hold, and let
\(\mathfrak B\) be a continuous box of \(\rho\) with good
residue.  Define \(Z_C^2\) by local constancy and the two
equations (APD.27).  Its coboundaries are the pairs
\((d\beta,a)\), where \(\beta\) is locally constant and
there are \(Y_v\) with
\[
 a_v(s)-\bigl(\beta(\operatorname{loc}_v s)
       -(\overline{\operatorname{Ad}}(\operatorname{loc}_v s)-1)Y_v\bigr)
                                                    \in\mathfrak b_v.
                                                               \tag{APD.37}
\]
Put \(H_C^2=Z_C^2/(Z_C^2\cap B_C^2)\), exactly the original
\(\mathrm{coneH2}\).  Then
\(\mathrm{HasObs}(\mathcal O,R,H_C^2)\) holds in the full
sense of APD5.  Its vanishing-to-section assertion is also the
full original \texttt{pr\_obs}.

\begin{proof}
For every original small surjection \(p:C'\to R\), choose
a set lift by APD14.  APD15 gives the linear map
\[
 \mathrm{ob}_p:K^\vee\longrightarrow H_C^2,
 \qquad u\longmapsto[(\omega_u,a_u)].                     \tag{APD.38}
\]
Suppose its value at \(u\) is zero, and let
\(C'\overset q\twoheadrightarrow C''\overset{p''}\to R\)
be any factorization permitted in APD5.  The map \(q|_K\)
factors through \(u\).  If \(u=0\), it is zero and we set
\(\mu=0\).  Otherwise \(u\) is surjective onto the
one-dimensional \(\kappa\)-space, and define
\(\mu(a)=q(y)\) whenever \(u(y)=a\).  The kernel condition
on \(q\) makes this independent of \(y\) and additive.
For every coefficient \(c''\in C''\), lift it by \(q\)
and reduce that lift modulo \(\mathfrak m_{C'}\).
The equality \(\mathfrak m_{C'}K=0\) then gives
\(\mu(\overline{c''}a)=c''\mu(a)\).
Its image lies in \(\mathfrak m_{C''}\), is contained in
\(q(K)\), and is killed by \(p''\).  This proves every
clause of the original \texttt{pr\_mu} construction.
Indeed the construction only uses an ideal
\(K\subset\mathfrak m_{C'}\) annihilated by
\(\mathfrak m_{C'}\), the surjection \(q\), and its stated
condition on \(\ker u\).  Without mentioning \(p,p''\), it
returns \(\mu(u(y))=q(y)\) and
\(\mu(\kappa)\subset q(K)\) at the full original general-ideal
scope of \texttt{pr\_mu}.

Push the chosen set lift through \(q\).  Its defects are
\(\mu(\omega_u)\) and \(\mu(a_u)\), because all original
defect entries are in \(K\).  Vanishing in the stated quotient
means there are actual \(\beta,Y\) satisfying (APD.37), with
\(\beta\) locally constant.  APD16 corrects the pushed set
lift to a continuous box.  Furthermore
\(\ker p''=q(K)\) and
\(\mathfrak m_{C''}\ker p''=0\), using surjectivity of \(q\).
APD17 therefore normalizes this box to a continuous
\(\rho''\) with \(p''_*\rho''=\rho\).

Residue preservation makes \(\bar\rho''=\bar\rho\).
Universality gives \(s:R\to C''\) with \(s_*\rho=\rho''\).
Both \(p''s\) and the identity of \(R\) induce \(\rho\),
so its uniqueness clause gives \(p''s=1\).  This is exactly
the section required in APD5, for every permitted quotient.
Nothing requires \(\ker q=\ker u\); the original implication
\(u(y)=0\Rightarrow q(y)=0\) is sufficient.

The argument from a supplied set lift and a supplied boundary
does not use surjectivity of \(p\).  Keeping that general map,
and only the stated surjectivity of \(q\), proves the original
\texttt{pr\_obs} as well.  In particular the proof establishes
the original obstruction predicate, not just liftability across
one chosen one-dimensional extension.
\end{proof}
% END APD18

% BEGIN APD19
\paragraph{APD19. Every strict tangent cocycle comes from the universal ring.}
\label{claim:APD19}
At the hypotheses of APD18, for every strict one-cocycle
\((z,(a_v))\in Z_C^1\) there are
\(D\in\operatorname{Der}_{\mathcal O}(R,\kappa)\) and
\(X\in M_n(\kappa)\) such that
\[
 \phi(D)(s):=D(b(s))\bar b(s)^{-1}
              =z(s)+(\overline{\operatorname{Ad}}(s)-1)X
                         \quad(s\in\Gamma).             \tag{APD.39}
\]
Before the universality step, the dual-number construction below
has the broader original \texttt{pr\_tbox} scope: it applies to
any box over \(\kappa\) with locally constant matrix part,
under the original frame and coefficient premises.  It needs no
\(\mathrm{Univ}\), \(\mathrm{Good}\), or
\(\mathrm{Box.Of}\) assumption.  Those additional hypotheses
enter only when the constructed box is normalized and represented
by an algebra map from \(R\).

The derivative acts entrywise on the matrix part.  This is full
original \texttt{pr\_tsurj}, not a surjectivity assumption on
the tangent map.

\begin{proof}
The residual box is locally constant.  Indeed its entries modulo
\(\mathfrak m_R\) are locally constant by continuity, and
there are finitely many matrix entries; the sign is locally
constant too.  Over \(\kappa[\epsilon]/\epsilon^2\) define
\[
 b_z(s)=(1+\epsilon z(s))\bar b(s),
 \qquad g_{z,v}=\bar g_v(1+\epsilon a_v).                  \tag{APD.40}
\]
The one-cocycle identity makes \(b_z\) multiplicative.  Its
local conjugate is
\((1+\epsilon\bar g_v[z(\operatorname{loc}_v s)
 -(\overline{\operatorname{Ad}}(\operatorname{loc}_v s)-1)a_v]
 \bar g_v^{-1})\bar t_{v,s}\), which is upper triangular
by (APD.11).  Thus (APD.40) is a continuous box reducing to
\(\bar{\mathfrak B}\).  It is precisely the original
\texttt{pr\_tbox} construction.

Apply APD17 to the reduction onto \(\kappa\).  Its kernel is
\((\epsilon)\), annihilated by the maximal ideal.  We obtain
a conjugating matrix \(h=1+\epsilon Y\) and a normalized
\(\rho_z\) with residue \(\bar\rho\).  Universality gives
an actual algebra map \(R\to\kappa[\epsilon]/\epsilon^2\)
inducing \(\rho_z\); APD1 expresses it as
\(\operatorname{res}_R+\epsilon D\).

Both the image of the original box and the conjugated box have
matrix \(1\) at \(c\) and agree on \(G_0\).  The decomposition
\(\Gamma=G_0\sqcup G_0c\) therefore makes their signed
representations equal on all of \(\Gamma\).  The first-order
coefficient of conjugation by \(1+\epsilon Y\) is
\(z(s)+Y-\overline{\operatorname{Ad}}(s)Y\).
Multiplying the matrix derivative by \(\bar b(s)^{-1}\)
therefore gives (APD.39) with \(X=-Y\).  This proves the
claimed tangent statement at arbitrary local \(\mathcal O\)
and at every rank satisfying the original good-residue premise.
\end{proof}
% END APD19

% BEGIN APD20
\paragraph{APD20. The weight torus supplies the extra tangent direction.}
\label{claim:APD20}
Let
\(T_0=\operatorname{diag}(w(n)_i)\in M_n(\kappa)\).
At the hypotheses of APD18, if
\(\theta_{\mathrm{Lie}}(T_0)=T_0\), there is a tangent
vector \(D_0\) with
\[
       \phi(D_0)(s)=(\overline{\operatorname{Ad}}(s)-1)T_0
                         \quad(s\in\Gamma).             \tag{APD.41}
\]
This is the full original \texttt{pr\_tor}.

\begin{proof}
Extend the residual box to the dual numbers and conjugate it by
\(h=1-\epsilon T_0\), changing every local basis to
\(\bar g_vh^{-1}\).  Its first-order cochain is exactly the
right side of (APD.41).  The condition
\(\theta_{\mathrm{Lie}}(T_0)=T_0\) ensures
\(\theta(h)=h\), so its matrix at \(c\) remains \(1\).
The marked difference is diagonal and \(h\) is diagonal;
thus this difference remains diagonal too.  APD13 consequently
recovers a normalized continuous representation from this box
without further normalization.  Its residual representation is
\(\bar\rho\).  Universality and APD1 give a tangent vector
\(D_0\), and equality of the two boxes on \(G_0\) and \(c\)
gives (APD.41) on \(\Gamma\), exactly as in APD19.

For the actual ranks, \(T_0=\operatorname{diag}(1,-1)\) or
\(\operatorname{diag}(1,0,-1)\).  It is nonzero over every
field, and reversing its diagonal entries and negating returns
itself.  These are the two properties needed by (APD.13).
Notice that \(D_0\) itself is allowed to be zero when
\(T_0\) is invariant; APD6 retains that branch and obtains
the extra dimension from \(H_C^0\) instead.
\end{proof}
% END APD20

\subsection{The residual calculation and the exact numerical difference}

% BEGIN APD21
\paragraph{APD21. The actual residue is good and has the required ordinary bases.}
\label{claim:APD21}
Return to the actual \(O',R,\rho\) of (APD.1), and write
\(k'=\operatorname{Res}(O')\).  There is the specified residue
map \(f:k\hookrightarrow k'\) such that
\(\bar\rho.\mathrm{IsBar}(f)\); the matrices of \(\bar\rho\)
are all upper triangular or all lower triangular.  Its residue is
good in the exact sense of APD17.  There is a continuous box
\(\mathfrak B\) of \(\rho\).  These statements supply full
\texttt{pr\_res}, the actual box input, and full original
\texttt{pr\_good} at every field extension of \(k\).

\begin{proof}
Take the residue map induced by the original local map \(O\to O'\).
ALC17 gives the residual diagonal and character formulas after
passing further to \(\operatorname{Res}(R)\).  The canonical
map \(k'\to\operatorname{Res}(R)\) is an isomorphism, so
injectivity descends the same formulas to \(k'\).  APD8 gives
the stated global triangular alternative at this exact residue.

Here is the good-residue calculation over any target field \(K\)
and any \(f:k\to K\).  The odd prime condition gives
\(2\ne0\).  If \(z=\Psi(g_0)\), then \(z^2\ne1\),
and the diagonal entries of the marked difference are
\(q,-q\), or \(q,0,-q\), where \(q=f(z-z^{-1})\ne0\).
They are pairwise distinct, since \(2q\ne0\).
For each local place the frame supplies \(s_v\) with
\(z_v=\Psi(\operatorname{loc}_v s_v)\) satisfying
\(z_v^2\ne1\).  The ordered character values at \(s_v\)
are \(f(z_v),f(z_v^{-1})\), or
\(f(z_v),1,f(z_v^{-1})\).  They are pairwise distinct:
equality of the outer pair gives \(z_v^2=1\), and equality
with \(1\) gives \(z_v=1\).  This proves all three
\(\mathrm{Good}\) fields.  The last two conclusions of
\texttt{pr\_good}, namely \(T_0\ne0\) and
\(\theta_{\mathrm{Lie}}T_0=T_0\), were explicitly checked
in APD20.

In \(R\), the differences of the full character values at
\(s_v\) are units, by the residual calculation.  The exact
characteristic-polynomial identity over \(\mathcal A\) is
ACP4 and survives \(\iota\).  Together with the ordered
ordinary-product identity it gives the ordered ordinary basis
over \(R\) by ACP3.  Continuity of \(\rho\) is ALC17.
APD13 constructs the required box from these bases.  Thus all
inputs here are supplied for the actual local representation;
no ordinary flag is left as a new hypothesis on it.
\end{proof}
% END APD21

% BEGIN APD22
\paragraph{APD22. The local and infinite adjoint spaces have their exact dimensions.}
\label{claim:APD22}
For every field \(K\), every rank, and every
\(g\in\mathrm{GL}_n(K)\), the subspace
\(\mathfrak b(g)=\{X:gXg^{-1}\text{ is upper triangular}\}\)
has dimension \(n(n+1)/2\).  For a box of a good normalized
representation, its local subspaces further satisfy
\[
 \dim_K\mathfrak b_v=\frac{n(n+1)}2,
 \qquad
 (\operatorname{Ad}(\operatorname{loc}_v s)-1)Y\in\mathfrak b_v
 \text{ for all }s\quad\Longrightarrow\quad Y\in\mathfrak b_v.
                                                               \tag{APD.42}
\]
The rigidity assertion requires only an arbitrary frame, a box
of the supplied normalized representation, and its
\(\mathrm{Good}\) condition; it does not require
\(F.\mathrm{Ok}\).  They are stable under the local adjoint
action.  This stability
requires only a box over any commutative ring and its defining
identities; neither good residue nor \(F.\mathrm{Ok}\) is
needed for that assertion.  If
\(x\in\mathrm{GG}_n(K)\), \(x^2=1\), its sign is negative,
and \(2\ne0\) in \(K\), then
\[
            \dim_K\ker(\operatorname{Ad}(x)-1)
                              =\frac{n(n-1)}2.            \tag{APD.43}
\]
For any continuous box over an \(\mathrm{IsCO}\)-algebra its
residual signed representation and its residual adjoint
representation have open kernels.  These assertions prove full
\texttt{pr\_bdim}, \texttt{pr\_brig}, \texttt{pr\_bstab},
\texttt{pr\_cfix}, and \texttt{pr\_open} at their respective
original scopes.

\begin{proof}
Conjugation by \(g_v\) identifies \(\mathfrak b_v\) with
the upper triangular matrices.  The entries on and above the
diagonal are free coordinates, giving \(n(n+1)/2\) dimensions
over any field.  Products and inverses of invertible upper
triangular matrices are upper triangular, so conjugation by a
local box matrix preserves this subspace, over any commutative
ring.

For the second assertion of (APD.42), take a separating local
element supplied by \(\mathrm{Good}\), and in the ordinary
basis write its matrix as \(A\) and the matrix of \(Y\) as
\(Z\).  Then \(A\) is upper triangular with distinct
diagonal entries.  The assumed condition, multiplied on the
right by \(A\), says that \(AZ-ZA\) is upper triangular.
For \(i>j\) its \((i,j)\) entry equals
\[
 (A_{ii}-A_{jj})Z_{ij}
       +\sum_{q>i}A_{iq}Z_{qj}-\sum_{q<j}Z_{iq}A_{qj}.
\]
Starting at the greatest distance below the diagonal and
descending, the summands in the two sums already vanish.
The distinct diagonal entries force \(Z_{ij}=0\).  This
proves the rigidity assertion for every rank, not only ranks
two and three.

For (APD.43), write the matrix part of \(x\) as \(A\).
The signed involution identity says \(A\theta(A)=1\),
equivalently \(S=AJ_n\) is an invertible symmetric matrix.
The adjoint action is \(X\mapsto-SX^tS^{-1}\).
Thus its fixed matrices are exactly those with
\((XS)^t=-XS\).  Right multiplication by \(S\) identifies
this space with all skew-symmetric matrices.  Since \(2\ne0\),
their diagonal entries are zero and each pair \(i<j\) supplies
one free coordinate.  There are \(n(n-1)/2\) such pairs.
This proves the original integer identity
\(2\dim\ker(\operatorname{Ad}(x)-1)+n=n^2\), including
the ranks zero and one.

Finally each residual matrix entry of a continuous box is locally
constant.  Their finite tuple is locally constant.  The kernel
of the signed representation is the intersection of its
identity-matrix fiber with the open kernel of \(\delta\),
so is open.  The adjoint kernel contains this open subgroup
and is a union of its cosets, so is open as well.  This also
proves the original local-constancy output, even when the
coefficient residue field is infinite.
\end{proof}
% END APD22

% BEGIN APD23
\paragraph{APD23. The cyclotomic dual of the residual adjoint module is zero.}
\label{claim:APD23}
Let \(n=2\) or \(3\), let \(K\) be any field of
characteristic \(\ell\), let \(f:k\to K\), and let
\(\rho\) satisfy \(\mathrm{IsBar}(f)\), with all its
matrices upper triangular or all lower triangular.  Let
\(\mathfrak B\) be a box of \(\rho\).  Every linear
functional \(\lambda:M_n(K)\to K\) satisfying
\[
 \lambda(\operatorname{Ad}(s)X)
                   =\omega(s)\lambda(X)
                         \quad(s\in\Gamma,X\in M_n(K)) \tag{APD.44}
\]
is zero.  Here \(\omega=F.\mathrm{om}(K)\) is the exact
cyclotomic character.  This is full original \texttt{pr\_dual}.

\begin{proof}
Write \(a_{ij}=\lambda(E_{ij})\).  Then
\(\lambda(X)=\sum_{i,j}X_{ij}a_{ij}\).  For \(g\in G_0\),
the adjoint formula and (APD.44), applied after multiplying by
\(\rho(g)\), give
\[
 \sum_p\rho(g)_{pi}a_{pj}
                  =\omega(g)\sum_q a_{iq}\rho(g)_{jq}.
                                                               \tag{APD.45}
\]
Put \(M(g)=\rho(g)^t\).  If these matrices are upper
triangular, induct on \(i\) from the last index downwards
and on \(j\) from the first index upwards.  Once the entries
already reached by that induction have vanished, (APD.45)
reduces to
\[
       \rho(g)_{ii}a_{ij}
                         =\omega(g)a_{ij}\rho(g)_{jj}
                                  \quad(g\in G_0).       \tag{APD.46}
\]
Indeed terms with \(p<i\) or \(q>j\) vanish by triangularity,
and the other off-diagonal terms are exactly the earlier ones.
If \(M\) is lower triangular, reverse both orders.  It remains
to show that (APD.46) forces \(a_{ij}=0\).

If \(a_{ij}\ne0\), substitute the prescribed diagonal
characters in (APD.46) and cancel this nonzero scalar.  The
cyclotomic character commutes with the field map from \(k\),
since it is induced from the prime field.  Injectivity of \(f\)
therefore gives
\[
          \omega_k(g)\Psi(g)^{w_j-w_i}=1
                                  \quad(g\in G_0).       \tag{APD.47}
\]
The possible differences are \(0,\pm1,\pm2\).  For zero,
the frame's \(f2a\) contradicts (APD.47); for \(\pm1\),
its \(f2b\) does.  If the difference is \(\pm2\), then
\(j=i^{\mathrm{rev}}\).  The matrix part of the box at
\(c\) is \(1\), and
\(\theta_{\mathrm{Lie}}(E_{i,i^{\mathrm{rev}}})
             =-E_{i,i^{\mathrm{rev}}}\).
Equation (APD.44) hence gives \(\omega(c)=-1\), because
the coefficient is nonzero.  With \(\Psi(c)=1\) and
\(\delta(c)=-1\), it follows that the character
\(\omega_k\Psi^{w_j-w_i}\delta_k\) is \(1\) both on
\(G_0\), by (APD.47), and on \(c\).  These generate
\(\Gamma\); it is identically \(1\), contradicting the
frame's \(f2c\).  This settles every weight difference.

The triangular induction now makes every \(a_{ij}\) zero,
so \(\lambda=0\).  In particular the two outer weights in
rank three have not been treated as adjacent weights: their
contradiction uses the negative-sign element and the precise
\(\delta_k\) factor.  No exclusion of \(\ell=3\) occurs.
\end{proof}
% END APD23

% BEGIN APD24
\paragraph{APD24. The tangent excess is at least \(nF.d+1\).}
\label{claim:APD24}
For the actual local ring and box supplied above, let
\[
 T=\operatorname{Der}_{O'}(R,k'),\quad
 h_i=\dim_{k'}H_C^i\ (i=0,1,2),\quad t=\dim_{k'}T.
\]
The obstruction space \(H_C^2\) is finite dimensional and
\[
                         h_2+nF.d+1\le t.                \tag{APD.48}
\]

\begin{proof}
All input scopes for the frame's Euler field can now be checked.
The field \(k'\) is finite of characteristic \(\ell\), by
the original coefficient-ring fields of \(L\).  The adjoint
space \(W=M_n(k')\) is finite dimensional.  Its kernel is
open and the local Borels are stable by APD22.  The residual
representation is globally triangular by APD21, and APD23
proves that its cyclotomic dual is zero.  Thus the literal
\(F.\mathrm{EC}\) field gives finite dimensionality of
\(H_C^2\) and
\[
 h_2+h_0+\sum_v F.dv(v)\dim\mathfrak b_v
        \le h_1+\sum_{w\in F.\mathrm{Inf}}
             \dim\ker(\overline{\operatorname{Ad}}(F.cw(w))-1).
                                                               \tag{APD.49}
\]
This is the previously reconstructed Euler field at exactly its
original finite-module premise, not a conjectural dimension formula.

Put \(a=n(n-1)/2\).  APD22 gives
\(\dim\mathfrak b_v=a+n\).  Each \(F.cw(w)\) is an
involution of negative sign, so its fixed adjoint dimension is
\(a\), again by APD22.  The original frame identities are
\(\sum_v F.dv(v)=F.d\) and
\(\#F.\mathrm{Inf}=F.d\).  Substitution in (APD.49)
and cancellation of \(F.d\,a\) yield
\[
                         h_2+h_0+nF.d\le h_1.             \tag{APD.50}
\]

APD12 supplies actual universality, APD21 supplies the good
residue, and APD1 makes \(T\) finite dimensional.  APD19
is the second premise of APD6; the local rigidity in APD22
is its first.  The nonzero \(T_0\) and the actual tangent
vector of APD20 supply its extra-direction premise.  Hence
\(h_1+1\le t+h_0\).  Adding one to (APD.50), combining
these inequalities, and cancelling \(h_0\) proves (APD.48).
This also displays exactly where the additional \(1\), beyond
\(nF.d\), enters the relation count.
\end{proof}
% END APD24

% BEGIN APD25
\paragraph{APD25. Full original \texttt{pres} on every supplied local object.}
\label{claim:APD25}
For every original \(O,F,C_0,n,S,x,L\) of (APD.1), there
are integers \(N,N'\), a surjective \(O'\)-algebra map
\[
 p:O'[[X_1,\ldots,X_N]]\longrightarrow R,
 \quad f_1,\ldots,f_{N'}\in O'[[X_1,\ldots,X_N]],
\]
such that
\[
            N=N'+nF.d+1,\qquad \ker p=(f_1,\ldots,f_{N'}).
                                                               \tag{APD.51}
\]
This is the full original \(L.\mathrm{Pres}\), with no
\(F.\mathrm{Big}\) or nontriangularity assumption.  In particular
it is available at every original rank-two and rank-three local
chart, including those arising from the permitted Host prime
\(\ell=3\).

\begin{proof}
APD18 supplies the original
\(\mathrm{HasObs}(O',R,H_C^2)\).  Set
\(q=nF.d+1\), \(t=\dim T\), and
\(h=\dim H_C^2\).  APD24 says \(h+q\le t\).
Let
\[
              H'=H_C^2\oplus(k')^{t-q-h}.
\]
For every small extension, compose the actual map (APD.38)
with the inclusion in the first summand of \(H'\).  This
inclusion is injective, so its zero locus is unchanged.
The same sections supplied in APD18 prove
\(\mathrm{HasObs}(O',R,H')\), with every quotient in its
original quantifier still allowed.

The ring \(O'\) is complete Noetherian local with finite
residue field.  APD5 therefore gives a surjection with exactly
\(t\) variables and \(\dim H'=t-q\) listed relations.
Take \(N=t\) and \(N'=t-q\).  Inequality (APD.48)
makes these natural-number differences legitimate, and
\(N=N'+q\) proves (APD.51).  If fewer nonzero generators
are necessary, the construction pads the family with zeros;
this preserves its literal generating-ideal assertion.  The
kernel is proper, and all these relations lie in the maximal
ideal of the power-series ring.

No replacement of the coefficient extension \(O'\), the
closed point \(x\), or the local map \(\iota\) has occurred.
The full ordinary characters entered point realization, the
box, normalization, and the obstruction corrections.  Thus
the result is a presentation of the actual original \(R\),
not of a deformation problem obtained by discarding characters
or by identifying only residual semisimplifications.
\end{proof}
% END APD25

\paragraph{The exact next consumers and the retained arithmetic frontier.}
The rank-two consumer \texttt{sharp2} in
\texttt{ChainedBot.lean} first extracts precisely (APD.51).
The power-series ring has dimension \(N+1\), so the next
component-chain input is used with \(h=N'\) and
\(D=2F.d+1\).  Its numerical threshold is
\(D+1+h=N+1\).  Exact numerical control of both counts
supplies this threshold.  The remaining original
\texttt{c1ch}, with its \texttt{ClComp}/\texttt{ClChain}
bodies, and the triangular and coincident-character dimension
suppliers remain separate work.  The present chapter does not
assert \texttt{sharp2}, \texttt{sharp3}, or a new component
dimension bound.

Rank three also uses the ordinary presentation in its appropriate
branch, but its even/odd \(\mathrm{EPres}\) record has additional
data and is not supplied by (APD.51).  Global closed-point and
seed incidence, the rank-two-to-rank-three corner, full
\(\mathrm{AllGood}\), actual \(\mathrm{OffFin}\), and the
automorphic propagation inputs remain visible.  The two spaces
and modularity ideals remain distinct on their common arithmetic
frame.  The arithmetic entry still retains the original \(P\),
and the Alt route still retains its own \(\mathrm{WitA}\) and
\(\mathrm{WitAQ}\) witnesses.  No whole selected terminal
obligation is closed by this local presentation chapter.
